\documentclass[12pt]{article}

\usepackage{sbc-template}
\usepackage{graphicx,url}
\usepackage[utf8]{inputenc}
\usepackage[T1]{fontenc}
\usepackage[brazil]{babel}
\usepackage{booktabs}
\usepackage{array}

\sloppy

\newcolumntype{L}[1]{>{\raggedright\arraybackslash}p{#1}}

% Citação narrativa no estilo SBC: "Autor et al. \citeano{chave}" -> "Autor et al. [2024]".
% Extrai o ano (última palavra) do rótulo do \bibitem.
\makeatletter
\def\sbc@lastword#1 #2\@nil{\if\relax\detokenize{#2}\relax#1\else\sbc@lastword#2\@nil\fi}
\newcommand{\citeano}[1]{\@ifundefined{b@#1}{[?]}{%
  [\expandafter\expandafter\expandafter\sbc@lastword\csname b@#1\endcsname{} \@nil]}%
  \nocite{#1}}
\makeatother

\title{Explicabilidade e Transparência em Agentes Baseados em LLM ao Longo
  do Ciclo de Vida de Desenvolvimento de Software (SDLC): Um Mapeamento
  Sistemático da Literatura}

\author{Davi Duarte Neco, Jordane Gabriella Soares de Oliveira Pires,\\
  Matheus Augusto Ferreira Medeiros, Matheus Vieira Mendes Pacheco}

% Cabeçalho só com título e autores, sem instituição nem e-mails.
\makeatletter
\def\@maketitle{\newpage
  \begin{center}
  \vspace*{-.7cm}
  {\XVIPT\bf\@title\par}
  \vglue 6pt plus 3pt minus 3pt
  {\normalsize\textbf{\begin{tabular}[t]{c}\@author\end{tabular}}\par}
  \vglue 6pt plus 3pt minus 3pt
  \end{center}\par}
\makeatother

\begin{document}

\maketitle

\begin{abstract}
  This paper presents a systematic mapping study on explainability and
transparency of LLM-based agents across the Software Development Life Cycle
(SDLC), covering the techniques employed, their distribution across SDLC phases,
their impact on developer trust and debugging, and how explanation quality is
evaluated. Searches in IEEE Xplore, ACM Digital Library and Google Scholar,
complemented by snowballing, yielded 17 primary studies (2022--2026). Agent
trajectory tracing and post-hoc attribution predominate, research concentrates
on implementation and maintenance, no study addresses requirements,
architecture or operation, and explanation faithfulness is rarely assessed.
\end{abstract}

\begin{resumo}
  Este artigo apresenta um mapeamento sistemático sobre explicabilidade e
transparência em agentes baseados em LLM ao longo do ciclo de vida de
desenvolvimento de software (SDLC), abrangendo as técnicas empregadas, sua
distribuição entre as fases do SDLC, seus impactos na confiança e na depuração e
a avaliação da qualidade das explicações. Buscas na IEEE Xplore, ACM Digital
Library e Google Scholar, complementadas por \textit{snowballing}, resultaram em
17 estudos primários (2022--2026). Predominam o rastreamento de trajetórias de
agentes e a atribuição pós-hoc, os estudos concentram-se em implementação e
manutenção, nenhum trata de requisitos, arquitetura ou operação, e a fidelidade
das explicações é pouco avaliada.
\end{resumo}

\section{Introdução}

A crescente adoção de \textit{Large Language Models} (LLMs) na Engenharia de
Software levou a agentes baseados em LLM que executam, com diferentes níveis de
autonomia, tarefas como análise de requisitos, geração e revisão de código,
testes, identificação de falhas e manutenção, atuando em diferentes etapas do
ciclo de vida de desenvolvimento de software (\textit{Software Development Life
Cycle}, SDLC).

O aumento dessa autonomia amplia a necessidade de compreender as decisões dos
agentes: não basta uma resposta aparentemente adequada, pois o engenheiro
precisa identificar quais informações fundamentaram uma decisão, quais etapas
foram executadas e o que contribuiu para o resultado. Isso torna relevantes a
transparência, a rastreabilidade e a explicabilidade, sobretudo porque LLMs
podem gerar informações incorretas ou não fundamentadas.

Embora existam revisões sobre o uso de LLMs na Engenharia de Software
\cite{hou2024,he2025} e sobre explicabilidade em modelos de aprendizado de
máquina aplicados à área \cite{mohammadkhani2023,cao2024}, não há um panorama que
relacione as técnicas de explicabilidade de agentes baseados em LLM às fases do
SDLC e sintetize como sua eficácia é avaliada.

Diante disso, este artigo apresenta um mapeamento sistemático da literatura,
conduzido segundo Petersen et al. \citeano{petersen2015} e Kitchenham e Charters
\citeano{kitchenham2007}, que analisa 17 estudos primários publicados entre 2022 e
2026. As contribuições são: (i) uma classificação das técnicas de
explicabilidade em cinco categorias; (ii) a distribuição dessas técnicas pelas
fases do SDLC; (iii) uma síntese das evidências sobre confiança, depuração e
excesso de confiança; e (iv) um levantamento dos métodos de avaliação das
explicações.

\section{Fundamentação teórica}\label{sec:fundamentacao}

\subsection{Agentes baseados em LLM na Engenharia de Software}

Um agente baseado em LLM é um sistema em que o modelo de linguagem planeja,
executa ações por meio de ferramentas (como edição de arquivos e execução de
testes) e observa seus resultados para decidir o passo seguinte \cite{yao2023}.
Na Engenharia de Software, há agentes capazes de resolver \textit{issues} reais
de repositórios, avaliados em \textit{benchmarks} como o SWE-bench, e revisões
recentes registram o uso de LLMs e de sistemas multiagentes em praticamente
todas as fases do SDLC \cite{hou2024,he2025}. Com menor supervisão humana,
cresce a preocupação com a confiança necessária para aceitar o que produzem
\cite{roychoudhury2025}. Neste mapeamento, consideram-se tanto agentes autônomos
quanto LLMs usados como assistentes, pois os mecanismos de explicação são, em
grande parte, os mesmos.

\subsection{Explicabilidade, transparência e interpretabilidade}

Os três termos têm significados distintos. Interpretabilidade é a capacidade de
apresentar o comportamento de um modelo em termos compreensíveis a um humano;
explicabilidade refere-se às explicações que justificam uma saída específica; e
transparência, à visibilidade sobre o funcionamento do sistema e o processo que
levou a um resultado \cite{barredo2020}. Em agentes, a transparência inclui a
rastreabilidade da trajetória: pensamentos, chamadas de ferramentas e
observações.

Em LLMs, as explicações incluem raciocínio verbalizado, como
\textit{chain-of-thought} (CoT) \cite{wei2022}; atribuição de importância à
entrada, como valores de Shapley \cite{lundberg2017}; proveniência das
informações em geração aumentada por recuperação (RAG) \cite{lewis2020}; e
sinais de incerteza. Uma questão central é a fidelidade (\textit{faithfulness}):
o quanto a explicação reflete o processo que de fato produziu a saída
\cite{jacovi2020}; explicações em CoT podem ser plausíveis e não fiéis
\cite{turpin2023}.

\subsection{Trabalhos relacionados}

Mohammadkhani et al. \citeano{mohammadkhani2023} revisaram 24 estudos sobre
explicabilidade na Engenharia de Software, concentrados em predição de
defeitos, e Cao et al. \citeano{cao2024} ampliaram o escopo para 108 estudos
baseados em aprendizado de máquina e profundo; ambas tratam sobretudo de modelos
anteriores aos LLMs e não consideram agentes. Já Hou et al. \citeano{hou2024} e He
et al. \citeano{he2025} mapeiam o uso de LLMs e de sistemas multiagentes, mas
abordam a explicabilidade apenas como desafio em aberto. Este trabalho situa-se
na interseção desses dois conjuntos.

\section{Metodologia}\label{sec:metodologia}

O estudo segue as diretrizes de Petersen et al. \citeano{petersen2015} para
mapeamentos sistemáticos, complementadas por Kitchenham e Charters
\citeano{kitchenham2007} na definição do protocolo. Seu objetivo é mapear as
abordagens de explicabilidade e transparência aplicadas a agentes baseados em
LLM na Engenharia de Software, identificando as técnicas utilizadas, sua
aplicação nas fases do SDLC e os métodos empregados para avaliar sua eficácia e
utilidade para desenvolvedores. Os registros de todas as etapas estão
disponíveis na planilha de seleção do repositório do estudo.

\subsection{Questões de pesquisa}

\begin{itemize}
  \item \textbf{RQ1 (Taxonomia e técnicas de XAI):} Quais são as técnicas de
    explicabilidade (ex.: raciocínio verbalizado/CoT, proveniência de RAG,
    árvores de decisão de agentes, saliência de código/\textit{tokens}) mais
    utilizadas para tornar transparentes as decisões de agentes LLM?
  \item \textbf{RQ2 (Mapeamento ao longo do SDLC):} Como as abordagens de
    explicabilidade variam entre as diferentes fases do ciclo de vida do
    software (ex.: explicação da escolha de arquitetura vs.\ justificativa para
    a geração de casos de teste ou correções de código)?
  \item \textbf{RQ3 (Eficácia e confiança do desenvolvedor):} Qual é o impacto
    documentado da explicabilidade dos agentes na tomada de decisão humana, na
    facilidade de depuração e na mitigação do excesso de confiança
    (\textit{over-reliance}) do engenheiro no código gerado?
  \item \textbf{RQ4 (Métricas e métodos de avaliação de XAI):} Como os estudos
    da literatura avaliam empiricamente a qualidade, a fidelidade e a utilidade
    das explicações fornecidas pelos agentes aos engenheiros de software (ex.:
    métricas automáticas, testes de uso com programadores, auditoria humana)?
\end{itemize}

\subsection{Estratégia de busca}

A \textit{string} de busca combina três grupos de termos derivados das questões
de pesquisa: (i) a tecnologia (LLMs e agentes de LLM); (ii) o atributo de
interesse (explicabilidade, transparência e interpretabilidade); e (iii) o
domínio (Engenharia e desenvolvimento de software). Os sinônimos de cada grupo
foram combinados por OR e os grupos, por AND:

\begin{quote}\small\ttfamily
("large language model" OR "LLM agent" OR "LLM agents") AND
(explainability OR transparency OR interpretability) AND
("software engineering" OR "software development")
\end{quote}

A \textit{string} foi aplicada a todos os campos, sem restrição de período;
termos de técnicas específicas (ex.: \textit{chain-of-thought}) foram omitidos
para não enviesar a busca. Uma busca exploratória em 1º/10/2026 calibrou a
\textit{string}; a busca definitiva, relatada aqui, foi executada em 15/11/2026
e complementada por \textit{snowballing} para trás e para frente
\cite{wohlin2014} a partir dos estudos incluídos.

\subsection{Fontes de busca}

Foram utilizadas a IEEE Xplore e a ACM Digital Library, que indexam os
principais veículos da área, e o Google Scholar, que abrange também
\textit{preprints}. As bases retornaram 65 (IEEE Xplore), 3.567 (ACM Digital
Library) e cerca de 17.700 (Google Scholar) resultados. Como os resultados são
ordenados por relevância e o volume é elevado, a triagem considerou os 8, 9 e 8
primeiros resultados de cada base, respectivamente.

\subsection{Critérios de inclusão e exclusão}

Foram incluídos os estudos que atendem a todos os critérios de inclusão (CI) e a
nenhum dos critérios de exclusão (CE):

\begin{itemize}
  \item[CI1 --] O estudo aborda LLMs ou agentes de LLMs.
  \item[CI2 --] O LLM/agente é aplicado a uma atividade de Engenharia de
    Software ou etapa do SDLC, como requisitos, testes, arquitetura,
    codificação, DevOps ou manutenção.
  \item[CI3 --] O estudo aborda algum mecanismo de explicabilidade,
    transparência, interpretabilidade, rastreabilidade ou proveniência
    relacionado ao comportamento/resultado do LLM ou agente.
  \item[CI4 --] O estudo apresenta informação suficiente para responder a pelo
    menos uma das RQs do mapeamento.
  \item[CI5 --] O texto completo está disponível para análise.
  \item[CI6 --] O estudo está escrito em inglês ou português.
\end{itemize}

\begin{itemize}
  \item[CE1 --] Estudo duplicado entre as bases.
  \item[CE2 --] Utiliza LLM, mas fora do contexto de Engenharia de Software.
  \item[CE3 --] Trata de Engenharia de Software e LLM, mas não aborda
    explicabilidade/transparência.
  \item[CE4 --] Trata de XAI/explicabilidade, mas não envolve LLMs ou agentes
    baseados em LLMs.
  \item[CE5 --] O LLM é usado apenas como instrumento auxiliar da pesquisa, sem
    que sua explicabilidade/transparência seja objeto do estudo.
  \item[CE6 --] Discute requisitos de explicabilidade de outros sistemas de IA,
    sem investigar a explicabilidade/transparência do LLM/agente atuando no
    SDLC.
  \item[CE7 --] Não possui texto completo acessível ou não apresenta
    informações suficientes para análise.
\end{itemize}

\subsection{Processo de seleção dos estudos}

A seleção teve quatro etapas: (1) reunião dos resultados e remoção de
duplicados (CE1); (2) aplicação dos critérios a título, resumo e
palavras-chave, encaminhando os casos duvidosos; (3) reaplicação dos critérios
aos casos encaminhados, com atenção a CE5 e CE6, que separam os estudos em que a
explicabilidade do LLM é objeto de investigação daqueles em que o LLM é apenas
instrumento; e (4) \textit{snowballing}. Cada decisão foi registrada com o
critério e uma justificativa.

A Tabela~\ref{tab:selecao} resume o fluxo. Na etapa 2, 14 dos 18 excluídos o
foram por CE3, isto é, a maioria dos resultados trata de LLMs em Engenharia de
Software sem abordar a explicabilidade. Na etapa 3, Dawoodjee et al.
\citeano{dawoodjee2026} e Mohammad et al. \citeano{mohammad2026} foram excluídos por
CE7, pois seus textos completos não estavam acessíveis.

O \textit{snowballing} partiu dos 5 estudos incluídos e foi iterado até que
nenhuma nova inclusão ocorresse: para trás, as listas de referências foram
triadas por título; para frente, foram verificadas as citações aos estudos
incluídos. Dos 19 estudos avaliados, 12 foram incluídos e 7 excluídos (4 por
CE4, por tratarem de modelos anteriores aos LLMs, e 1 por CE1, CE2 e CE3 cada).

\begin{table}[ht]
\centering
\caption{Fluxo da seleção dos estudos}\label{tab:selecao}
\small
\begin{tabular}{lcccc}
\toprule
Etapa & IEEE & ACM & Scholar & Total \\
\midrule
Resultados triados (título e resumo) & 8 & 9 & 8 & 25 \\
Excluídos na etapa 2 & 7 & 8 & 3 & 18 \\
Encaminhados à etapa 3 & 1 & 1 & 5 & 7 \\
Excluídos na etapa 3 (CE7) & 1 & 0 & 1 & 2 \\
Incluídos a partir das bases & 0 & 1 & 4 & 5 \\
\midrule
\multicolumn{4}{l}{Avaliados no \textit{snowballing}} & 19 \\
\multicolumn{4}{l}{Incluídos via \textit{snowballing}} & 12 \\
\multicolumn{4}{l}{\textbf{Total de estudos incluídos}} & \textbf{17} \\
\bottomrule
\end{tabular}
\end{table}

\subsection{Extração e classificação dos dados}

Para cada estudo, registraram-se: identificação; nível de autonomia (assistente
ou agente); técnica (RQ1); fase do SDLC (RQ2); efeitos sobre confiança, decisão,
depuração e excesso de confiança (RQ3); e método de avaliação (RQ4).

As técnicas foram classificadas por codificação aberta seguida de agrupamento,
partindo dos exemplos da RQ1, o que resultou em cinco categorias: T1 --
justificativa em linguagem natural (raciocínio verbalizado, CoT ou explicação
textual da saída); T2 -- atribuição pós-hoc (importância de \textit{tokens},
trechos ou características da entrada); T3 -- sinalização de incerteza e
confiança; T4 -- rastreamento de trajetórias e proveniência (registro e análise
dos passos do agente ou da origem do artefato); e T5 -- elicitação de
necessidades de explicação. Cada estudo foi associado à técnica predominante. As
fases do SDLC seguem Hou et al. \citeano{hou2024}, acrescidas de revisão de código
e de uma categoria transversal.

\section{Resultados}\label{sec:resultados}

Os 17 estudos incluídos (Tabela~\ref{tab:estudos}) foram publicados entre 2022 e
2026 (1 em 2022, 2 em 2024, 9 em 2025 e 5 em 2026); onze passaram por revisão
por pares e seis são \textit{preprints}.

\begin{table}[ht]
\centering
\caption{Estudos primários incluídos no mapeamento}\label{tab:estudos}
\scriptsize
\begin{tabular}{@{}lL{2.4cm}L{3.4cm}cL{1.2cm}L{3.9cm}@{}}
\toprule
ID & Estudo & Fase / atividade & Téc. & Origem & Avaliação das explicações \\
\midrule
E01 & Sun et al. \citeano{sun2022} & Implementação (geração) & T5 & SB & Workshops com 43 engenheiros \\
E02 & Palacio et al. \citeano{palacio2024}$^\dagger$ & Implementação (geração) & T2 & SB & Questionário (24); análise causal \\
E03 & Li et al. \citeano{li2024} & Manutenção (sumarização) & T2 & SB & \textit{Eye-tracking} vs.\ atribuição \\
E04 & Virk et al. \citeano{virk2025} & Manutenção (sumarização) & T3 & SB & Calibração (Brier, ECE) \\
E05 & Kang et al. \citeano{kang2025} & Depuração e reparo & T1 & Scholar & Estudo com 20 participantes \\
E06 & Vasconcelos et al. \citeano{vasconcelos2025} & Implementação (completação) & T3 & SB & Experimento controlado (30) \\
E07 & Spiess et al. \citeano{spiess2025} & Implementação (geração) & T3 & SB & Calibração (várias tarefas) \\
E08 & Wang et al. \citeano{wang2025} & Implementação (geração) & T4 & ACM & Sem avaliação (proposta) \\
E09 & Ceka et al. \citeano{ceka2025}$^\dagger$ & Manutenção (\textit{issues}) & T4 & SB & Trajetórias de 5 agentes \\
E10 & Bouzenia e Pradel \citeano{bouzenia2025} & Manutenção (\textit{issues}) & T4 & SB & 120 trajetórias anotadas \\
E11 & Majgaonkar et al. \citeano{majgaonkar2025}$^\dagger$ & Manutenção (\textit{issues}) & T4 & SB & Trajetórias de 3 agentes \\
E12 & Khati et al. \citeano{khati2025}$^\dagger$ & Implementação e testes & T2 & SB & Entropia; estudo com usuários (37) \\
E13 & Vitale et al. \citeano{vitale2026} & Implementação (geração) & T2 & Scholar & Fidelidade; questionário (37) \\
E14 & Akbarova et al. \citeano{akbarova2026} & Testes (fronteira) & T1 & SB & Questionário (27); entrevistas (6) \\
E15 & Joshi \citeano{joshi2026}$^\dagger$ & Depuração (falhas de agentes) & T4 & Scholar & Estudo com 20 participantes \\
E16 & Li e Storhaug \citeano{listorhaug2026} & Transversal (avaliação de agentes) & T4 & Scholar & 18 artigos; estudo de caso \\
E17 & Gao et al. \citeano{gao2026}$^\dagger$ & Revisão de código & T1 & SB & Experimento \textit{within-subjects} (34) \\
\bottomrule
\multicolumn{6}{@{}L{14.2cm}@{}}{$^\dagger$ \textit{Preprint} sem revisão por pares até a data da busca. SB: \textit{snowballing}. T1: justificativa em linguagem natural; T2: atribuição pós-hoc; T3: incerteza e confiança; T4: trajetórias e proveniência; T5: necessidades de explicação.}
\end{tabular}
\end{table}

\subsection{RQ1: Técnicas de explicabilidade}

O rastreamento de trajetórias e proveniência (T4) é a categoria mais frequente
(6 estudos). Quatro estudos analisam sequências de pensamento, ação e resultado
de agentes que resolvem \textit{issues}
\cite{bouzenia2025,ceka2025,majgaonkar2025} ou falham em tarefas de programação
\cite{joshi2026}, convertendo, neste último, os registros brutos em explicações
estruturadas. Li e Storhaug \citeano{listorhaug2026} defendem publicar essas
trajetórias para explicar por que uma abordagem agêntica supera outra, e Wang et
al. \citeano{wang2025} propõem embutir a rastreabilidade entre requisito e código no
\textit{pipeline} de geração.

A atribuição pós-hoc (T2, 4 estudos) inclui valores de Shapley sobre
características da entrada \cite{vitale2026}, confiança por \textit{token}
alinhada à árvore sintática \cite{palacio2024}, agregação global de explicações
por \textit{token} \cite{khati2025} e a comparação entre os trechos considerados
importantes pelo modelo e pelos programadores \cite{li2024}. A justificativa em
linguagem natural (T1, 3 estudos) aparece na depuração por hipótese e
experimento \cite{kang2025}, em testes de fronteira \cite{akbarova2026} e na
revisão de código \cite{gao2026}. A sinalização de incerteza (T3, 3 estudos)
compreende o destaque de \textit{tokens} propensos a edição
\cite{vasconcelos2025} e a calibração da confiança de resumos \cite{virk2025} e
de código \cite{spiess2025}. Por fim, Sun et al. \citeano{sun2022} (T5) elicitaram
necessidades de explicação com 43 engenheiros.

Das técnicas citadas na RQ1, árvores de decisão de agentes não foram
encontradas, e a proveniência em RAG aparece apenas como proposta, sem
avaliação \cite{wang2025}. Além disso, apenas 6 dos 17 estudos (E05, E09, E10,
E11, E15 e E16) tratam de agentes com autonomia de múltiplos passos; os demais
investigam LLMs usados como assistentes.

\subsection{RQ2: Explicabilidade ao longo do SDLC}

A Tabela~\ref{tab:fases} relaciona as técnicas às fases do SDLC. Implementação
e manutenção concentram 12 dos 17 estudos, e nenhum trata de requisitos,
arquitetura ou operação.

\begin{table}[ht]
\centering
\caption{Distribuição dos estudos por técnica e fase do SDLC (n = 17)}\label{tab:fases}
\footnotesize\setlength{\tabcolsep}{4pt}
\begin{tabular}{lccccccc}
\toprule
Técnica & Implem. & Manut. & Depur. & Testes & Revisão & Transv. & Total \\
\midrule
T1 -- Justificativa em linguagem natural & -- & -- & 1 & 1 & 1 & -- & 3 \\
T2 -- Atribuição pós-hoc & 3 & 1 & -- & -- & -- & -- & 4 \\
T3 -- Incerteza e confiança & 2 & 1 & -- & -- & -- & -- & 3 \\
T4 -- Trajetórias e proveniência & 1 & 3 & 1 & -- & -- & 1 & 6 \\
T5 -- Necessidades de explicação & 1 & -- & -- & -- & -- & -- & 1 \\
\midrule
Total & 7 & 5 & 2 & 1 & 1 & 1 & 17 \\
\bottomrule
\end{tabular}
\end{table}

A técnica empregada varia com a fase. A atribuição pós-hoc (T2) e a sinalização
de incerteza (T3) concentram-se na implementação e na sumarização de código, em
que a saída é um artefato que pode ser ancorado em trechos da entrada. O
rastreamento de trajetórias (T4) predomina na manutenção, associado a agentes
que resolvem \textit{issues} em vários passos. Nas atividades em que a saída é
um julgamento (depuração, testes e revisão de código), predomina a
justificativa em linguagem natural (T1): explica-se por que um \textit{patch}
corrige a falha \cite{kang2025}, por que um par de entrada e saída é uma
fronteira \cite{akbarova2026} ou por que um trecho deve ser alterado
\cite{gao2026}.

\subsection{RQ3: Impacto na confiança, na decisão e na depuração}

Os estudos com desenvolvedores indicam efeitos positivos, mas não uniformes. Na
depuração, as explicações do AutoSD aumentaram a acurácia ao julgar a corretude
de \textit{patches} em cinco de seis \textit{bugs}, sem aumentar o tempo, e 70\%
dos 20 participantes quiseram explicações \cite{kang2025}.
Explicações estruturadas a partir de trajetórias de agentes permitiram
identificar a causa-raiz de falhas 2,8 vezes mais rápido e propor correções com
acurácia 73\% maior do que com os registros brutos \cite{joshi2026}. Na revisão
de código (34 participantes), explicações completas geraram a maior confiança
percebida, mas a maior concordância com a IA ocorreu com explicações moderadas,
sugerindo que explicações detalhadas levam a questionar mais \cite{gao2026}. Na
completação de código, destacar os \textit{tokens} com maior probabilidade de
edição tornou as edições mais direcionadas e as tarefas mais rápidas, ao
contrário de destacar \textit{tokens} de baixa probabilidade de geração
\cite{vasconcelos2025}.

Nem todos os resultados são favoráveis: o foco atribuído ao modelo na
sumarização de código não corresponde ao foco dos programadores medido por
\textit{eye-tracking} \cite{li2024}. Quanto ao excesso de confiança, as
evidências são exploratórias, mas apontam riscos: participantes editaram apenas
os trechos destacados e ignoraram um caso de borda, lendo a ausência de destaque
como sinal de corretude \cite{vasconcelos2025}; uma única explicação alucinada
reduziu fortemente as avaliações dos participantes \cite{akbarova2026}; e os
modelos de código não são bem calibrados sem ajuste, o que limita o uso direto
de sua confiança como sinal \cite{spiess2025}. Nenhum estudo mediu o excesso de
confiança de forma longitudinal ou em contexto real de desenvolvimento.

\subsection{RQ4: Métodos de avaliação das explicações}

Os estudos combinam três tipos de avaliação (Tabela~\ref{tab:estudos}): estudos
com desenvolvedores (10), métricas automáticas (6) e análise qualitativa de
trajetórias (4); apenas Wang et al. \citeano{wang2025} não apresentam avaliação.
Nos estudos com usuários predominam questionários Likert sobre clareza,
utilidade e confiança; experimentos controlados com medidas de desempenho são
menos frequentes \cite{kang2025,vasconcelos2025,gao2026,joshi2026}, e as amostras
variam de 20 a 43 participantes. As métricas automáticas incluem calibração
\cite{virk2025,spiess2025}, entropia \cite{khati2025}, fidelidade
\cite{vitale2026}, inferência causal \cite{palacio2024} e o alinhamento entre a
atribuição do modelo e o foco humano \cite{li2024}.

Apenas dois estudos avaliaram explicitamente a fidelidade das explicações, ambos
em T2: Vitale et al. \citeano{vitale2026}, por perturbação da entrada, e Palacio et
al. \citeano{palacio2024}, por inferência causal. Nenhum estudo de T1 avaliou a
fidelidade das justificativas, embora explicações plausíveis possam não ser
fiéis \cite{turpin2023}. Nos estudos de T4, a trajetória é tratada como registro
fiel do processo, sem verificar se o raciocínio verbalizado corresponde às
ações executadas.

\section{Ameaças à validade}\label{sec:ameacas}

As ameaças foram analisadas segundo as categorias de Ampatzoglou et al.
\citeano{ampatzoglou2019}. Quanto à identificação dos estudos, a \textit{string}
não inclui variações como \textit{explainable}, \textit{explanation}, \textit{rationale} ou
\textit{trajectory}, e a triagem foi limitada aos primeiros resultados por
relevância, o que pode ter omitido estudos, como indica a origem de 12 dos 17
estudos no \textit{snowballing}, cuja direção para frente se restringiu às
citações entre os estudos identificados. Além disso, uma versão preliminar
deste trabalho incluía cinco estudos (sobre rastreabilidade de requisitos,
causa-raiz de incidentes, localização de falhas, \textit{design rationale} e
revisão de código) não alcançados pelo processo documentado, que foram
removidos; a ausência de requisitos, arquitetura e operação na
Tabela~\ref{tab:fases} deve ser lida à luz dessa limitação.

Quanto à seleção, a distinção entre CE5 e CE6 envolve julgamento, por isso cada
decisão foi justificada; dois estudos potencialmente relevantes foram excluídos
apenas por indisponibilidade do texto (CE7). Quanto à extração, as categorias
T1--T5 foram definidas pelos autores, e registrar só a técnica predominante pode
subestimar combinações (ex.: Joshi \citeano{joshi2026}). Os seis
\textit{preprints} podem mudar após a revisão por pares. Por fim, as
frequências relatadas são indicativas, não representativas.

\section{Considerações finais}\label{sec:conclusao}

Este artigo apresentou um mapeamento sistemático sobre explicabilidade e
transparência em agentes baseados em LLM ao longo do SDLC. Nos 17 estudos
analisados, predominam o rastreamento de trajetórias e a atribuição pós-hoc; as
pesquisas concentram-se em implementação e manutenção; as explicações tendem a
melhorar a acurácia do julgamento e a depuração, mas podem induzir dependência
dos sinais exibidos; e a avaliação apoia-se em percepções de usuários e métricas
automáticas, com pouca atenção à fidelidade.

Como lacunas, destacam-se a ausência de estudos sobre árvores de decisão de
agentes e proveniência em RAG, a escassez de explicações nas fases iniciais do
SDLC, a falta de avaliação da fidelidade das justificativas e trajetórias e a
inexistência de estudos longitudinais sobre excesso de confiança. Como trabalhos
futuros, pretende-se ampliar a \textit{string}, incluir a Scopus, realizar o
\textit{snowballing} para frente sobre um índice de citações e avaliar a
qualidade dos estudos primários.

\begin{thebibliography}{99}

\bibitem[Akbarova et al. 2026]{akbarova2026}
Akbarova, S., Dobslaw, F., and Feldt, R. (2026). Understanding on the edge:
LLM-generated boundary test explanations. In \textit{Proc. IEEE/ACM Int. Conf.
on Automation of Software Test (AST)}.

\bibitem[Ampatzoglou et al. 2019]{ampatzoglou2019}
Ampatzoglou, A. et al. (2019). Identifying, categorizing and mitigating threats
to validity in software engineering secondary studies. \textit{Information and
Software Technology}, 106:201--230.

\bibitem[Barredo Arrieta et al. 2020]{barredo2020}
Barredo Arrieta, A. et al. (2020). Explainable artificial intelligence (XAI):
Concepts, taxonomies, opportunities and challenges toward responsible AI.
\textit{Information Fusion}, 58:82--115.

\bibitem[Bouzenia and Pradel 2025]{bouzenia2025}
Bouzenia, I. and Pradel, M. (2025). Understanding software engineering agents: A
study of thought-action-result trajectories. In \textit{Proc. 40th IEEE/ACM Int.
Conf. on Automated Software Engineering (ASE)}.

\bibitem[Cao et al. 2024]{cao2024}
Cao, S. et al. (2024). A systematic literature review on explainability for
machine/deep learning-based software engineering research. arXiv:2401.14617.

\bibitem[Ceka et al. 2025]{ceka2025}
Ceka, I. et al. (2025). Understanding software engineering agents through the
lens of traceability: An empirical study. arXiv:2506.08311.

\bibitem[Dawoodjee et al. 2026]{dawoodjee2026}
Dawoodjee, I. E., Khurat, A., and Noraset, T. (2026). A multimodal large language
model approach for automated security requirements verification. In \textit{Proc.
23rd Int. Joint Conf. on Computer Science and Software Engineering (JCSSE)}.

\bibitem[Gao et al. 2026]{gao2026}
Gao, Z. et al. (2026). Evaluating the impact of explainable AI on trust in
AI-assisted code review. arXiv:2607.24601.

\bibitem[He et al. 2025]{he2025}
He, J., Treude, C., and Lo, D. (2025). LLM-based multi-agent systems for software
engineering: Literature review, vision, and the road ahead. \textit{ACM Trans.
Softw. Eng. Methodol.}, 34(5).

\bibitem[Hou et al. 2024]{hou2024}
Hou, X. et al. (2024). Large language models for software engineering: A
systematic literature review. \textit{ACM Trans. Softw. Eng. Methodol.}, 33(8).

\bibitem[Jacovi and Goldberg 2020]{jacovi2020}
Jacovi, A. and Goldberg, Y. (2020). Towards faithfully interpretable NLP systems:
How should we define and evaluate faithfulness? In \textit{Proc. 58th Annual
Meeting of the ACL}, pages 4198--4205.

\bibitem[Joshi 2026]{joshi2026}
Joshi, A. (2026). XAI for coding agent failures: Transforming raw execution
traces into actionable insights. arXiv:2603.05941.

\bibitem[Kang et al. 2025]{kang2025}
Kang, S. et al. (2025). Explainable automated debugging via large language
model-driven scientific debugging. \textit{Empirical Software Engineering},
30(2):45.

\bibitem[Khati et al. 2025]{khati2025}
Khati, D. et al. (2025). Enabling global, human-centered explanations for LLMs:
From tokens to interpretable code and test generation. arXiv:2503.16771.

\bibitem[Kitchenham and Charters 2007]{kitchenham2007}
Kitchenham, B. and Charters, S. (2007). Guidelines for performing systematic
literature reviews in software engineering. Tech. Report EBSE-2007-01, Keele
University and Durham University.

\bibitem[Lewis et al. 2020]{lewis2020}
Lewis, P. et al. (2020). Retrieval-augmented generation for knowledge-intensive
NLP tasks. In \textit{Advances in Neural Information Processing Systems
(NeurIPS)}, 33.

\bibitem[Li et al. 2024]{li2024}
Li, J. et al. (2024). Do machines and humans focus on similar code? Exploring
explainability of large language models in code summarization. In \textit{Proc.
32nd IEEE/ACM Int. Conf. on Program Comprehension (ICPC)}.

\bibitem[Li and Storhaug 2026]{listorhaug2026}
Li, J. and Storhaug, A. (2026). Reproducible, explainable, and effective
evaluations of agentic AI for software engineering. In \textit{Companion Proc.
34th ACM Int. Conf. on the Foundations of Software Engineering (FSE)}.

\bibitem[Lundberg and Lee 2017]{lundberg2017}
Lundberg, S. M. and Lee, S.-I. (2017). A unified approach to interpreting model
predictions. In \textit{Advances in Neural Information Processing Systems
(NeurIPS)}, 30.

\bibitem[Majgaonkar et al. 2025]{majgaonkar2025}
Majgaonkar, O. et al. (2025). Understanding code agent behaviour: An empirical
study of success and failure trajectories. arXiv:2511.00197.

\bibitem[Mohammad et al. 2026]{mohammad2026}
Mohammad, F. et al. (2026). CodeQual-Agent: An intelligent LLM-agent framework
for automated software quality assessment with explainable predictions and
real-time analysis. In \textit{IEEE Int. Conf. on Software Testing, Verification
and Validation Workshops (ICSTW)}.

\bibitem[Mohammadkhani et al. 2023]{mohammadkhani2023}
Mohammadkhani, A. H. et al. (2023). A systematic literature review of explainable
AI for software engineering. arXiv:2302.06065.

\bibitem[Palacio et al. 2024]{palacio2024}
Palacio, D. N. et al. (2024). Towards more trustworthy and interpretable LLMs for
code through syntax-grounded explanations. arXiv:2407.08983.

\bibitem[Petersen et al. 2015]{petersen2015}
Petersen, K., Vakkalanka, S., and Kuzniarz, L. (2015). Guidelines for conducting
systematic mapping studies in software engineering: An update.
\textit{Information and Software Technology}, 64:1--18.

\bibitem[Roychoudhury et al. 2025]{roychoudhury2025}
Roychoudhury, A. et al. (2025). Agentic AI software engineers: Programming with
trust. \textit{Communications of the ACM}.

\bibitem[Spiess et al. 2025]{spiess2025}
Spiess, C. et al. (2025). Calibration and correctness of language models for
code. In \textit{Proc. 47th IEEE/ACM Int. Conf. on Software Engineering (ICSE)}.

\bibitem[Sun et al. 2022]{sun2022}
Sun, J. et al. (2022). Investigating explainability of generative AI for code
through scenario-based design. In \textit{Proc. 27th Int. Conf. on Intelligent
User Interfaces (IUI)}, pages 212--228.

\bibitem[Turpin et al. 2023]{turpin2023}
Turpin, M. et al. (2023). Language models don't always say what they think:
Unfaithful explanations in chain-of-thought prompting. In \textit{Advances in
Neural Information Processing Systems (NeurIPS)}, 36.

\bibitem[Vasconcelos et al. 2025]{vasconcelos2025}
Vasconcelos, H. et al. (2025). Generation probabilities are not enough:
Uncertainty highlighting in AI code completions. \textit{ACM Trans.
Comput.-Hum. Interact.}, 32(1).

\bibitem[Virk et al. 2025]{virk2025}
Virk, Y., Devanbu, P., and Ahmed, T. (2025). Calibration of large language models
on code summarization. \textit{Proc. ACM on Software Engineering}, 2(FSE).

\bibitem[Vitale et al. 2026]{vitale2026}
Vitale, A. et al. (2026). Toward explaining large language models in software
engineering tasks. \textit{ACM Trans. Softw. Eng. Methodol.}

\bibitem[Wang et al. 2025]{wang2025}
Wang, F. et al. (2025). Embedding traceability in large language model code
generation: Towards trustworthy AI-augmented software engineering. In
\textit{Companion Proc. 33rd ACM Int. Conf. on the Foundations of Software
Engineering (FSE)}, pages 1760--1763.

\bibitem[Wei et al. 2022]{wei2022}
Wei, J. et al. (2022). Chain-of-thought prompting elicits reasoning in large
language models. In \textit{Advances in Neural Information Processing Systems
(NeurIPS)}, 35.

\bibitem[Wohlin 2014]{wohlin2014}
Wohlin, C. (2014). Guidelines for snowballing in systematic literature studies
and a replication in software engineering. In \textit{Proc. 18th Int. Conf. on
Evaluation and Assessment in Software Engineering (EASE)}.

\bibitem[Yao et al. 2023]{yao2023}
Yao, S. et al. (2023). ReAct: Synergizing reasoning and acting in language
models. In \textit{Int. Conf. on Learning Representations (ICLR)}.

\end{thebibliography}

\end{document}
